\paragrafos{Πρόσθεση ρητών: κέρδη, ζημιές και ισορροπία}{Α.7.3, «Πρόσθεση ρητών αριθμών»}

\begin{anagke}
Στο πούλμαν η παρέα παίζει ένα παιχνίδι με κάρτες: κάθε κάρτα γράφει έναν πόντο, θετικό για κέρδος και αρνητικό για ζημιά. Ο Βασίλης τραβάει $+5$ και μετά $-3$. «Πόσους πόντους έχω;» ρωτάει. Ο Νίκος υπολογίζει: «Πέντε μείον τρία κάνει δύο». Η Πολυξένη τραβάει $-4$ και μετά $-2$ και ο Στέλιος λέει «έξι», μα η Μυρτώ τον ρωτάει: «Έξι κέρδος ή έξι ζημιά;». Ο Ηλίας κάνει μια πιο δύσκολη ερώτηση: «Αν τραβήξω $-7$ και μετά $+3$, είμαι ακόμα πίσω;».

Η Αγγελική το βλέπει στον άξονα: «Κάθε κάρτα είναι μια μετακίνηση. Το $+5$ σημαίνει «πέντε βήματα δεξιά», το $-3$ «τρία βήματα αριστερά». Ξεκινάω από το μηδέν, πηγαίνω πέντε δεξιά και μετά τρία αριστερά και καταλήγω στο $+2$». Για το $-4$ και το $-2$ είναι έξι βήματα αριστερά: $-6$. Και για τα $-7$ και $+3$ καταλήγουμε στο $-4$: ακόμα πίσω, αλλά λιγότερο.

Με αυτόν τον τρόπο η παρέα ανακαλύπτει τους κανόνες της πρόσθεσης ρητών και μαθαίνει γιατί ισχύουν. Μάλιστα στο τέλος του παιχνιδιού ο Νίκος παρατηρεί κάτι ωραίο: αν μια κάρτα $+3$ και μια $-3$ βρεθούν μαζί, «ακυρώνονται».
\end{anagke}

\begin{lexi}{πρόσθεση, προσθετέος, άθροισμα}{πρόσθεση, άθροισμα}
Η \lexiword{πρόσθεση} είναι το «προσ-τίθημι», δηλαδή «βάζω επιπλέον». Οι \lexiword{προσθετέοι} είναι αυτά που προσθέτουμε και το \lexiword{άθροισμα} είναι αυτό που βγαίνει όταν τα «αθροίζουμε», τα βάζουμε σε ένα σωρό. Στους ρητούς, το «προσθέτω» μπορεί να σημαίνει και «ζημιώνω».
\end{lexi}

\begin{metafrasi}
$5$ βήματα δεξιά και $3$ αριστερά. & $(+5)+(-3)=+2$\\
$4$ βήματα αριστερά και $2$ αριστερά ακόμα. & $(-4)+(-2)=-6$\\
Κέρδος και ζημιά που ακυρώνονται. & $(+3)+(-3)=0$\\
Δεν αλλάζει η σειρά των προσθετέων. & $\alpha+\beta=\beta+\alpha$\\
Το $0$ δεν αλλάζει τίποτα. & $\alpha+0=\alpha$\\
\end{metafrasi}

\tmima{Η θεωρία}

Στον άξονα, προσθέτω έναν θετικό όταν πηγαίνω προς τα δεξιά και έναν αρνητικό όταν πηγαίνω προς τα αριστερά, ξεκινώντας από τη θέση του πρώτου προσθετέου. Από αυτό προκύπτουν οι κανόνες:

\begin{kanonas}{πρόσθεση ρητών}
\begin{itemize}[leftmargin=1.2em,itemsep=0.3ex,topsep=2pt]
\item \textbf{Ομόσημοι:} προσθέτω τις απόλυτες τιμές και κρατάω το κοινό πρόσημο. $(-8)+(-5)=-(8+5)=-13$ και $(+9)+(+4)=+13$.
\item \textbf{Ετερόσημοι:} αφαιρώ τη μικρότερη απόλυτη τιμή από τη μεγαλύτερη και βάζω το πρόσημο του προσθετέου με τη μεγαλύτερη απόλυτη τιμή. $(+9)+(-4)=+5$ και $(-9)+(+4)=-5$.
\item \textbf{Αντίθετοι:} το άθροισμα είναι $0$. $(-3{,}5)+(+3{,}5)=0$.
\item Με το $0$: $\alpha+0=\alpha$.
\end{itemize}
\end{kanonas}

\begin{idiotita}{ιδιότητες της πρόσθεσης}
Για όλους τους ρητούς $\alpha,\,\beta,\,\gamma$ ισχύουν: \textbf{αντιμεταθετική} $\alpha+\beta=\beta+\alpha$, \textbf{προσεταιριστική} $(\alpha+\beta)+\gamma=\alpha+(\beta+\gamma)$, $\alpha+0=\alpha$ και $\alpha+(-\alpha)=0$.
\end{idiotita}

Χάρη στις ιδιότητες, όταν έχω πολλούς προσθετέους μπορώ να τους ομαδοποιήσω όπως θέλω: προσθέτω χωριστά όλους τους θετικούς και χωριστά όλους τους αρνητικούς και μετά τα δύο αποτελέσματα. Έτσι $(-7)+12+(-5)+9+(-3)=(12+9)+\big((-7)+(-5)+(-3)\big)=21+(-15)=6$.

Η πρόσθεση κλασμάτων και δεκαδικών ακολουθεί τους ίδιους κανόνες: $\left(-\frac{1}{2}\right)+\left(-\frac{1}{4}\right)=-\frac{3}{4}$ και $\left(+\frac{2}{3}\right)+\left(-\frac{1}{2}\right)=\frac{4}{6}-\frac{3}{6}=\frac{1}{6}$.

\begin{prosoxi}
\begin{itemize}[leftmargin=1.2em,itemsep=0.3ex,topsep=2pt]
\item Το άθροισμα δύο αριθμών δεν είναι πάντα μεγαλύτερο από τους προσθετέους: $(-3)+(-2)=-5<-3$.
\item Στους ετερόσημους δεν προσθέτω τις απόλυτες τιμές: $(+9)+(-4)\neq13$, αλλά $9-4=5$.
\item Το πρόσημο του αθροίσματος ετεροσήμων είναι του μεγαλύτερου κατ' απόλυτη τιμή: $(-9)+4=-5$, όχι $+5$.
\item Στο $5+(-3)$ οι παρενθέσεις δείχνουν ότι το $-3$ είναι ο αριθμός που προσθέτω, όχι δύο πράξεις.
\end{itemize}
\end{prosoxi}

\tmima{Λυμένα παραδείγματα}

\paradeigma{Πρόσθεση στον άξονα}
\tixero Οι προσθέσεις $(-2)+5$, $3+(-7)$, $(-4)+(-3)$.\par
\tithelo Να τις κάνω στον άξονα.\par
\possindeo Ξεκινάω από τον πρώτο προσθετέο και κινούμαι: από το $-2$ πέντε δεξιά: $3$. Από το $3$ επτά αριστερά: $-4$. Από το $-4$ τρία αριστερά: $-7$.\par
\apantisi $3$, $-4$, $-7$.\par
\elegxos Με τους κανόνες: $5-2=3$, $-(7-3)=-4$, $-(4+3)=-7$.

\paradeigma{Με τους κανόνες}
\tixero Οι αριθμοί $(-8)+(-5)$, $(+9)+(-4)$, $(-9)+(+4)$, $\left(-\frac{1}{2}\right)+\left(-\frac{1}{4}\right)$.\par
\tithelo Να υπολογίσω τα αθροίσματα.\par
\possindeo Ομόσημοι: $-(8+5)=-13$. Ετερόσημοι με μεγαλύτερη απόλυτη τιμή το $+9$: $+(9-4)=5$. Ετερόσημοι με μεγαλύτερη το $-9$: $-(9-4)=-5$. Κλάσματα ομόσημα: $-\left(\frac{2}{4}+\frac{1}{4}\right)=-\frac{3}{4}$.\par
\apantisi $-13$, $5$, $-5$, $-\frac{3}{4}$.\par
\elegxos Κάθε πρόσημο το βρίσκω από τον προσθετέο με τη μεγαλύτερη απόλυτη τιμή.

\paradeigma{Πολλοί προσθετέοι}
\tixero Το άθροισμα $(-7)+12+(-5)+9+(-3)$.\par
\tithelo Να το υπολογίσω.\par
\possindeo Ομαδοποιώ: θετικοί $12+9=21$. Αρνητικοί $(-7)+(-5)+(-3)=-15$. Τότε $21+(-15)=6$.\par
\apantisi $6$.\par
\elegxos Με τη σειρά: $-7+12=5$, $5-5=0$, $0+9=9$, $9-3=6$.

\paradeigma{Οι ιδιότητες βοηθούν}
\tixero Το άθροισμα $(-4{,}6)+(+3{,}9)+(+4{,}6)+(-3{,}9)$.\par
\tithelo Να το βρω εύκολα.\par
\possindeo Με την αντιμεταθετική και την προσεταιριστική ομαδοποιώ τους αντίθετους: $\big((-4{,}6)+(+4{,}6)\big)+\big((+3{,}9)+(-3{,}9)\big)=0+0=0$.\par
\apantisi $0$.\par
\elegxos Ο καθένας έχει τον αντίθετό του και ακυρώνεται.

\paradeigma{Το ταμείο της παρέας}
\tixero Το ταμείο έχει $+30$ €. Κινήσεις: $-12$, $+8$, $-35$, $+20$.\par
\tithelo Το υπόλοιπο.\par
\possindeo Θετικοί: $30+8+20=58$. Αρνητικοί: $-12-35=-47$. $58+(-47)=11$.\par
\apantisi $11$ €.\par
\elegxos Βήμα βήμα: $30,\ 18,\ 26,\ -9,\ 11$. Συμφωνεί.

\tmima{Ασκήσεις}

\zoni{Το κατάλαβα}{κανόνες, ιδιότητες, πολλοί προσθετέοι}

\ex Στον άξονα βρες τα αθροίσματα: α) $(-2)+5$ \quad β) $3+(-7)$ \quad γ) $(-4)+(-3)$ \quad δ) $6+(-6)$ \quad ε) $(-1)+(+1)$.

\ex Υπολόγισε: α) $(-8)+(-5)$ \quad β) $(+9)+(-4)$ \quad γ) $(-9)+(+4)$ \quad δ) $(+7)+(+5)$ \quad ε) $(-12)+0$ \quad στ) $(-2{,}5)+(+1{,}5)$.

\ex Υπολόγισε: α) $(-3{,}5)+(+3{,}5)$ \quad β) $(-0{,}7)+(-1{,}8)$ \quad γ) $(+2{,}4)+(-5{,}6)$ \quad δ) $\left(-\frac{1}{2}\right)+\left(-\frac{1}{4}\right)$ \quad ε) $\left(+\frac{2}{3}\right)+\left(-\frac{1}{2}\right)$ \quad στ) $\left(-\frac{5}{6}\right)+\left(+\frac{1}{3}\right)$.

\ex Συμπλήρωσε τον αριθμό που λείπει ώστε το άθροισμα να είναι $0$: α) $(-7)+\kouti$ \quad β) $(+3{,}5)+\kouti$ \quad γ) $\kouti+(-4)$ \quad δ) $\left(-\frac{1}{3}\right)+\kouti$.

\ex Υπολόγισε: α) $(-7)+12+(-5)+9+(-3)$ \quad β) $(-1{,}5)+4+(-2{,}5)+1$ \quad γ) $10+(-3)+(-4)+(-5)$.

\ex Υπολόγισε με τις ιδιότητες: α) $(-4{,}6)+(+3{,}9)+(+4{,}6)+(-3{,}9)$ \quad β) $37+(-18)+(-37)+8$.

\ex Σωστό ή λάθος; α) Το άθροισμα δύο αρνητικών είναι αρνητικό. \quad β) Το άθροισμα δύο ετεροσήμων είναι πάντα αρνητικό. \quad γ) $(-5)+5=0$. \quad δ) $(-3)+(-4)=+7$. \quad ε) $\alpha+0=\alpha$. \quad στ) Το άθροισμα δύο αντίθετων είναι $0$. \quad ζ) $(-2)+(+5)>0$.

\zoni{Το χρησιμοποιώ}{θερμόμετρα, ταμεία, παιχνίδια}

\ex \textbf{Η θερμοκρασία.} Το πρωί στο Μέτσοβο είναι $-6$°, το μεσημέρι ανεβαίνει $9$ βαθμούς και το βράδυ πέφτει $5$ βαθμούς. Πόση είναι η θερμοκρασία το μεσημέρι και το βράδυ;

\ex \textbf{Το ταμείο.} Το ταμείο έχει $+30$ €. Γίνονται οι κινήσεις $-12$, $+8$, $-35$, $+20$. Πόσο είναι το υπόλοιπο μετά από κάθε κίνηση και στο τέλος;

\ex \textbf{Ο ανελκυστήρας.} Ο Νίκος είναι στον όροφο $-2$, ανεβαίνει $5$ ορόφους και μετά κατεβαίνει $4$. Σε ποιον όροφο φτάνει;

\ex \textbf{Ο δύτης.} Ο δύτης είναι στα $-18$ μ., ανεβαίνει $7$ μ. και μετά κατεβαίνει $4$ μ. Πού βρίσκεται;

\ex \textbf{Το παιχνίδι.} Οι πόντοι της Πολυξένης στους τέσσερις γύρους είναι $+8$, $-5$, $-7$, $+12$. Πόσους έχει συνολικά;

\ex \textbf{Ο λογαριασμός.} Η Μυρτώ έχει υπόλοιπο $-45$ € στον λογαριασμό, καταθέτει $80$ € και μετά αναλαμβάνει $20$ €. Πόσα έχει;

\ex \textbf{Η διαφορά καλαθιών.} Η ομάδα μπάσκετ έχει στους πέντε αγώνες διαφορά $+12$, $-7$, $-9$, $+4$, $+6$. Ποια είναι η συνολική διαφορά;

\zoni{Το σκέφτομαι}{γιατί ισχύουν οι κανόνες}

\ex Πότε το άθροισμα δύο ρητών είναι $0$;

\ex Αν οι $\alpha,\,\beta$ είναι ετερόσημοι, ποιο πρόσημο έχει το άθροισμά τους;

\ex Γράψε δύο διαφορετικά ζεύγη ρητών με άθροισμα $-5$.

\ex Είναι πάντα το άθροισμα μεγαλύτερο από κάθε προσθετέο; Δώσε αντιπαράδειγμα.

\ex Δείξε με τα $\alpha=3$ και $\beta=4$ ότι $(-\alpha)+(-\beta)=-(\alpha+\beta)$.

\zoni{Για το σπίτι}{ο προϋπολογισμός σου}

\ex Κράτησε για μια εβδομάδα τα έσοδα και έξοδα του χαρτζιλικιού σου ως ρητούς και υπολόγισε το τελικό υπόλοιπο.

\ex Φτιάξε ένα παιχνίδι με κάρτες $+$ και $-$ και παίξε με έναν φίλο τρεις γύρους. Ποιος έχει το μεγαλύτερο άθροισμα;
